\chapter{Tracing frameworks}

The term \textit{tracing framework} refers to a set of APIs and behaviours
designed to interface event sources with other kernel objects and user
front-ends. It is an infrastructural component because it provides the
architecture upon which tracing happens. In general, a tracing framework should
support the registration and de-registration of event sources, whether static or
dynamic, enabling and disabling, and producing or consuming their side effects.
Usually, user-space interacts with the functionalities exposed by a tracing
framework through a well-defined stable user-space API, such as through a
filesystem or a dedicated system call. Notably, there is no single, unified
tracing framework on Linux for registering and accessing all kinds of events,
but there are multiple solutions tailored to different use cases. The purpose of
this chapter is to provide an overview of the design and implementation of the
two major tracing frameworks present in Linux: \textit{ftrace} and
\textit{perf}.

\section{Ftrace}

\textit{Ftrace} is a versatile tracing framework that supports registration,
deregistration, reading, writing and other operations on various types of trace
events. It defines a \textit{tracer} abstraction to manage a set of events in
groups, and a \textit{global ftrace ring buffer} to produce and consume all the
event data. Additionally, it implements functionalities used by many types of
tracers and events, such as the registration of \textit{hooks} (altering a
funciton's execution to jump to a callback) on most kernel functions and other
code patching routines. Ftrace is most useful for tracing global events that
happen on any CPU and by any process, although it can be configured to filter
events with specific conditions, enabling ftrace to be adapted to various
use-cases.

Ftrace was originally implemented in the Linux $2.6$ kernel to support
\textit{function tracing}, that is, capturing a funciton's execution data such
as its arguments and name whenever the function is entered or exited, in real
time. Over time, the same infrastructure was reused for other use cases and
tracers, consequently the name "ftrace" does not retain any of its original
meaning. For this reason, the term "function tracer" will be used to refer to
the specific tracer used for functions, and "ftrace" for the wider
infrastructure.

\subsection{The ftrace global ring buffer}

A core component of ftrace is its \textit{global ring buffer}, managed at its
highest abstraction level by the \verb|struct trace_array| type. It provides an
interface to write data and consume it, supporting both a producer-consumer mode
and an overwrite mode. The two modes are distinguished by how writers behaves
when the buffer is full: in producer-consumer mode the producer fails to write,
otherwise it overwrites the data.

Since events are expected to happen at a high frequency and from many different
execution contexts, the design of the ring buffer needs to meet specific
requirements. These requirements include:

\begin{itemize}

  \item Writers must handle reentrancy safely, ensuring successful execution
    even when interrupted by hardware interrupts or non-maskable interrupts
    (NMIs).

  \item The design must support high concurrency for multiple writers and
    readers distributed across multiple CPU cores.

  \item Synchronization overhead must be minimized to efficiently handle a high
volume of concurrent writes from high-frequency events.

\end{itemize}

The solution implemented in Linux in the file
\verb|linux/kernel/trace/ring_buffer.c| fulfills all these requirements using a
lockless design for both reader-to-writer and writer-to-writer interactions.
Some key design decisions are briefly discussed.

Each CPU uses its own per-CPU ringbuffer containing a doubly-linked list of
pages; consequently, writers on different CPUs do not need to synchronize among
each other. In order to read a full trace in chronological order, a reader
iterates over all the per-CPU buffers using time information to determine the
order of the events. Another key detail is to store the reader page externally
to the ring buffer linked list: writers never interact with the reader page and
readers swap the reader page via atomic operations so no lock is used on the
same CPU ring buffer.

\subsubsection*{Reserve and commit pattern}

To support preemption, writers use a \textit{reserve and commit} pattern which
defines three phases to perform a write: first the writer reserves allocation
space in a per-CPU ring buffer, then it writes with exclusive access to that
region, and finally commits the region. The commit will be in a pending state
until all regions before it are also committed; therefore when a writer reserves
a region but gets preempted by another writer, then the new writer will not have
access to the reserved region and its changes will not get registered until the
original writer completes. This mechanism ensures that readers only process
fully committed entries.

In C code, the commit pattern looks like the following function from
\textit{\$LINUX/linux/kernel/trace/ trace\_uprobe.c}, commented and simplified
for readability:

\begin{Verbatim}[numbers=left, numbersep=5pt, frame=single]
static void __uprobe_trace_func(struct trace_uprobe *tu,
        unsigned long func, struct pt_regs *regs,
        struct uprobe_cpu_buffer *ucb,
        struct trace_event_file *trace_file)
{
  struct trace_event_buffer fbuffer;
  void *data;
  int size = ...;

  /* First, reserve an area */
  entry = trace_event_buffer_reserve(&fbuffer, trace_file, size);
  if (!entry)
    return;

  data = /* get the data pointer from and entry */

  /* Second, copy the data into the ring buffer page */
  memcpy(data, ucb->buf, ucb->dsize);

  /* Third, commit */
  trace_event_buffer_commit(&fbuffer);
}
\end{Verbatim}

The relevant interactions with the ring buffer are: the reservation on line
$11$, the writing on line $18$ via a \verb|memcpy()| and the commit on line
$21$.

\subsubsection*{Trace sequences}

Because string formatting is expensive, the event is stored in the ring buffer
in binary form without any formatting. When an end-user reads the ring buffer
though the user-space API, the relevant contents of the ring buffer are
formatted with a printf-style function, \verb|trace_seq_vprintf()|, into a trace
sequence of type \verb|struct trace_seq|. A trace sequence stores a pre
allocated buffer of 8K bytes that can be written to and read from. After
formatting, the data is copied from the trace sequence to user memory. Moving
the formatting procedure to the reader improves the writer's latency which is
critical for high frequency events.

\subsection{The tracer struct}

An important abstraction in ftrace is the \textit{tracer}, implemented as the
\verb|struct tracer| type. This object provides an interface to enable or
disable a set of trace events, with custom formatting if necessary. Any tracer
implements a set of functions to be invoked during its lifetime, such as
initialization, starting, stopping, reading, etc. An use case for a tracer would
be, for example, to trace all functions executed, or to trace all context
switches. The user interacts with the tracer via the \textit{tracefs} filesystem
API which will bee discussed later. Many tracers are available, but only one
tracer can be enabled at a time.

The default tracer is the \textit{nop} tracer which is a "dummy" tracer that does
not enable any event. The instantiation of such tracer looks like this:

\begin{Verbatim}[numbers=left, numbersep=5pt, frame=single]
/* Supported output options */
enum {
  TRACE_NOP_OPT_ACCEPT = 0x1,
  TRACE_NOP_OPT_REFUSE = 0x2
};

/* Options for the tracer */
static struct tracer_opt nop_opts[] = {
  { TRACER_OPT(test_nop_accept, TRACE_NOP_OPT_ACCEPT) },
  { TRACER_OPT(test_nop_refuse, TRACE_NOP_OPT_REFUSE) },
  { } /* Always set a last empty entry */
};

static struct tracer_flags nop_flags = {
  .val = 0, /* By default: all flags disabled */
  .opts = nop_opts
};

struct tracer nop_trace __read_mostly =
{
  .name = "nop",
  .init = nop_trace_init,  /* dummy functions, do not do anything */
  .reset = nop_trace_reset,
  .flags = &nop_flags,
  .set_flag = nop_set_flag,
  .allow_instances = true,
};
\end{Verbatim}

The \verb|nop_trace| tracer is instantiated in line $19$, where the \verb|init|
and \verb|reset| members on line $22$ and $23$ respectively are pointers to
functions that are supposed to bring up and remove all events managed by the
tracer, but, in this implementation, do nothing. The tracer is registered to the
ftrace infrastructure via the \verb|register_tracer()| function which is called
during kernel startup.

The following is a non exhaustive list of available tracers:

\begin{itemize}

\item \textit{function}: trace all kernel functions when entered.

\item \textit{function\_graph}: trace both function entry and exit. It can be
used to draw a graph of function calls.

\item \textit{blk}: block tracer, for detailed IO tracing.

\item \textit{hwlat}: hardware latency tracer.

\item \textit{irqsoff}: trace the areas that disable interrupts and saves the
trace with the longest maximum latency.

\item \textit{preemptoff}: similar to \textit{irqsoff} but traces and records
the amount of time for which preemption is disabled.

\item \textit{wakeup}: traces and records the maximum latency that it takes for
the highest priority task to get scheduled after it has been woken up.

\end{itemize}

\subsection{Static function tracing}

To support tracing function enter and exit events, ftrace implements a system to
instrument any function in order to alter its execution and trace the event. For
clarity, the following section will use the term \textit{trace} function or
\textit{tracing} function to refer to the function that performs the writing of
the event onto the ftrace ringbuffer, and the \textit{traced} function to the
function that is being instrumented (traced).

When function tracing is enabled, whenever a C function is executed such that
that the program counter reaches the first instruction belonging the function,
then the execution should jump to a tracing function and return to the traced
function without altering the registers. Likewise, when the program counter
reaches the end of the function to return, it should jump to a tracing function
and later jump to the original return address. This behaviour should be
supported transparently, without altering existing code or to expect developers
to statically place instrumentation onto any function.

This can be implemented thanks to special compiler support for this specific use
case. When \textit{gcc} is used with the \verb|-pg| command line flags, a jump
to a trace function called \verb|mcount()| will be generated after the prologue
of any function on the ARM CPU architecture, or a jump to \verb|__fentry__()|
before the prologue on the x86-64 architecture. For function exits, the trace
function will change the return address of the traced function with the address
of a trace function. This replacement can happen because the return address of a
function is saved on the stack as of C calling conventions, so changing this
stack value will alter execution after returning from the function. This is
referred to as "static function tracing" because the compiler instruments each
function permanently at compile time, and can be enabled via the
\verb|CONFIG_FUNCTION_TRACER| kernel compile option.

For example, the following \verb|sum()| function is compiled by gcc 16.1 into
x86-64 assembly with \verb|-O1| flag as shown below. On the left the code in C,
on the right its compiled code in assembly:

\noindent
\vspace{0.5em}

\begin{minipage}[t]{0.47\textwidth}
\begin{verbatim}
int sum(int a, int b) {
    return a + b;
}
\end{verbatim}
\end{minipage}%
\hfill
\vrule % This creates the vertical dividing bar
\hfill
\begin{minipage}[t]{0.47\textwidth}
\begin{verbatim}
"sum(int, int)":
        lea     eax, [rdi+rsi]
        ret
\end{verbatim}
\end{minipage}

\noindent
\vspace{1.0em}

The same compiler would compile the same function as follows, this time with the
\verb|-pg| flag added:

\vspace{1.0em}

\begin{minipage}[t]{0.47\textwidth}
\begin{verbatim}
int sum(int a, int b) {
    return a + b;
}
\end{verbatim}
\end{minipage}%
\hfill
\vrule % This creates the vertical dividing bar
\hfill
\begin{minipage}[t]{0.47\textwidth}
\begin{verbatim}
"sum(int, int)":
        call    __fentry__
        lea     eax, [rdi+rsi]
        ret
\end{verbatim}
\end{minipage}

\vspace{1.0em}

In the second version, the compiler added a call to \verb|__fentry__()| before
performing the sum operation. The \verb|__fentry__()| function can be
implemented anywhere and run any arbitrary code, which in the case of Linux
would be used for tracing purposes.

A shortcoming of this implementation is that each function suffers the cost of a
jump and a return at minimum, causing some performance loss in case tracing is
compiled in but not turned on. This issue is solved with "dynamic function
tracing", which will be discussed next.

\subsection{Dynamic function tracing}

In order to not have any performance penalty when tracing is supported but
disabled, ftrace provides a way to dynamically patch functions to add or remove
tracing instrumentation. Dynamic tracing is enabled with the kernel compile
option \verb|DYNAMIC_FTRACE|.

When dynamic function tracing is enabled, the kernel will patch all calls to
\verb|mcount()| or \verb|__fentry__()| with \verb|NOP| operations at boot-time,
which do not perform any operation. When tracing is enabled, the \verb|NOP|
instructions are patched with a call to \verb|ftrace_caller()| or
\verb|ftrace_regs_caller()| which can access a copy of all original registers
from the traced function. The information of all the locations of all
\verb|mcount()| calls is generated at compile time via the \verb|recordmcount|
program which is invoked by the build infrastructure. The tool stores all the
locations of these functions in the \verb|__mcount_loc| ELF section, which is
linked in the original object file. Through this information, ftrace know what
and where to patch code. This system allows for minimum overhead in the case
where tracing is disabled.

The tracing function \verb|ftrace_caller| is also referred to as a "trampoline"
because its instructions are dynamically patched to jump to the functions which
are currently enabled depending on the current tracer o particular events,
represented in the \verb|struct trace_ops|. Ftrace provides the function
\verb|register_ftrace_function()| to register a trampoline which is implemented
by an architecture dependent function, \verb|create_trampoline()|.

For example, consider the following assembly function:

\begin{verbatim}
some_kernel_func:
  PUSH RBP
  MOV RSP, RBP
  ...
  RET
\end{verbatim}

When tracing is disabled but dynamic tracing is supported, some \verb|NOP|
operations would be patched in at the start like in the next example:

\begin{verbatim}
# Symbol where instructions like: CALL some_kernel_func will jump to.
some_kernel_func:
  NOP
  NOP
# actual start of the function
  PUSH RBP
  MOV RSP, RBP
  ...
  RET
\end{verbatim}

When tracing is enabled, these same \verb|NOP| functions are replaced with a
call to \verb|ftrace_caller()|:

\begin{verbatim}
# Symbol where instructions like: CALL some_kernel_func will jump to.
some_kernel_func:
  CALL ftrace_caller
  NOP
# actual start of the function
  PUSH RBP
  MOV RSP, RBP
  [...]
  RET
\end{verbatim}

\subsection{The tracefs filesystem}

The ftrace subsystem needs to provide an interface for the user to interact with
it. This is done through the \textit{tracefs} filesystem. Compared to other user
interfaces like systemcalls and ioctls, many Linux kernel developers prefer
custom filesystems or custom entries in an already existing filesystem as a
method to provide new functionalities to the user. This is the case because
filesystems make the interface explicit and visible in user space; they also
make it easier to write scripts which perform administrative functions.
Originally, ftrace used the \textit{debufgs} filesystem but it grew enough to
motivate the addition of the current interface, tracefs. In this section, the
term \textit{\$TRACEFS} will be used to represent the location where the tracefs
filesystem is mounted.

Users can access the formatted content of all ring buffers by reading the
\textit{\$TRACEFS/trace} file without consuming the buffer, or
\textit{\$TRACEFS/trace\_pipe} to consume it. Additionally, they can get
information about the buffer state from several files beginning with
\verb|buffer_*|. To access a particular per-CPU buffer, they can use the files
from the directory \textit{\$TRACEFS/per\_cpu/\$CPU\_NUMBER/}. Users can query a
list of the tracers supported by the running kernel by reading the file
\textit{\$TRACEFS/available\_tracers}. They can select a tracer by writing its
name to the file \textit{\$TRAEFS/current\_tracer} and tune its options from
\textit{\$TRACEFS/options/\$OPTION\_NAME}. Additionally, tracefs provides a way
to specify filters for events, so that only certain evens in certain contexts
will be logged.

A more thorough explanation of the tracefs interface is available in the kernel
documentation.

\section{Perf}

The \textit{perf} subsystem is another tracing framework provided by the kernel
to manage and interact with trace events. It was originally designed as an
interface to read hardware performance counters (hence the name "perf" as in
"performance"), which are increasing integers provided by the hardware
manufacturers that store various information about the hardware such as number
of CPU cache misses or number of incorrect branch predictions. Counters are very
useful for performance analysis as they provide precise measurements. They are
simple because they only signal that an event has happened and how many times,
but the perf subsystem can also be used to interface with other types of events,
such as static or dynamic events that log rich information, therefore providing
a different interface than ftrace. Event counters as provided by perf are
referred to as \textit{perf events}, which can be either \textit{software} or
\textit{hardware} depending on what generated the event or increased the
counter.

\subsection{The Perf ring buffer}

To produce and consume event data, the perf infrastructure implements its own
ring buffer data structure, with some similarities and key differences compared
to ftrace's implementation. This is used for \textit{sampling} events that
periodically write measurements to the buffer, as opposed to \textit{counting}
events which count the aggregate number of events that occur and commonly do not
use the ring buffer. Unlike ftrace's, the perf ring buffer is not a global
static ring buffer, but it is rather created by the user via a systemcall and
liked to the lifetime of another object, such as a file descriptor. Another key
difference from ftrace's is the fact that the perf ring buffer is a contiguous
area of virtual memory that can be memory mapped via \verb|mmap()| to the user's
virtual address space, while ftrace uses a linked list of pages which are read
through a filesystem. Memory mapping the ring buffer removes the need of
repeated \verb|read()| systemcalls which are expensive operations. The memory
mapped area follows the structure of \verb|struct perf_event_mmap_page|
containing a pointer to the "head" producer, which is updated by the kernel, and
"tail" consumer, updated by the user program. Consequently, it is the
responsibility of the user to atomically update the tail pointer after consuming
the data, while the kernel updates the head pointer as it fills the buffer.
Similarly to ftrace, the perf ring buffer uses a reserve and commit pattern with
the \verb|perf_output_begin()| and \verb|perf_output_end()| functions.

For hardware events, perf uses architecture-specific \textit{auxiliary (AUX)
ring buffers}. Since they are implemented based on a particular architecture's
needs, they provide lower latency for high frequency events compared to a
generic implementation.

\subsection{The Performance Monitoring Unit}

To manage perf events, the perf subsystem introduces the \textit{Performance
Monitoring Unit (PMU)} abstraction implemented as the \verb|struct pmu| type.
This object manages the entire life cycle of an event, from registering and
unregistering it, to starting and stopping tracing, and reading the counter.

All the PMUs are registered with the \verb|perf_pmu_register()| function and
added to a global list of all PMUs, called "\verb|pmus|". The function takes a
\verb|struct pmu|, a name and an integer variable named "type". The name is used
to register a directory in the \textit{sysfs} filesystem under
\textit{\$SYSFS/bus/event\_source/devices/\$NAME/}, where \verb|$SYSFS| is the
location of the sysfs filesystem, and \verb|$NAME| is the name of the PMU passed
to the registration function. Through the sysfs API the user can enumerate and
get useful metadata about the events, but reading the actual event must be done
through a systemcall as described in the next section. The "type" argument will
specify which binary format the event will use. Each PMU contains pointer to
functions to manage one or more events; the structure looks like the following:

\begin{verbatim}
struct pmu {
  // ...
  void (*pmu_enable) (struct pmu *pmu); /* optional */
  void (*pmu_disable) (struct pmu *pmu); /* optional */
  int  (*event_init) (struct perf_event *event);
  int  (*add) (struct perf_event *event, int flags);
  void (*del) (struct perf_event *event, int flags);
  void (*start) (struct perf_event *event, int flags);
  void (*stop) (struct perf_event *event, int flags);
  // ...
};
\end{verbatim}

A more detailed look at how an event counter is registered will be taken in a
later chapter.

\subsection{The perf\_event\_open syscall}

The user interacts with the perf subsystem with a specific systemcall:
\verb|perf_event_open()|. Through this API, the user can access event counters
or to the ring buffer with control over which context should be traced such per
per-thread or per-CPU. This systemcall does not have a \textit{libc} wrapper.

The function signature of the systemcall is the following:

\begin{verbatim}
int perf_event_open(struct perf_event_attr *attr,
                    pid_t pid, int cpu, int group_fd,
                    unsigned long flags);
\end{verbatim}

The \verb|struct perf_event_attr| type is used to specify which event should be
monitored and some variables such as the frequency of sampling or whether to
exclude counting on certain contexts such as user or kernel. The full definition
of the structure is described in the man page. The "pid" and "cpu" arguments are
used to specify which process and CPU to monitor, the allowed values and their
meanings are also explained in the manual page. Finally, the "group\_fd"
argument can be used to create events in groups which is useful to measure
multiple events simultaneously, and the "flags" argument specify additional
options.

The return value is a new file descriptor or $-1$ on error. The user can use the
file descriptor to read the counter, or to memory map the ring buffer to
userspace to access its memory directly. Calls to \verb|ioctl()| using the file
descriptor can be issued to enable or disable the event. Multiple counters can
be kept open at a time, and the counters can be polled with \verb|poll()|.

\subsection{Security considerations}

Performance counters introduces substantial security risks regarding the
potential leakage of sensitive data provided by monitored events. These events
can leak information such as hardware configuration, content of registers,
execution data (context switches, pages faults) etc, with various degrees of
risk. Therefore, the perf subsystem supports several security frameworks: the
Unix privilege model, POSIX capabilities and Linux security modules (LSM).

In particular, usage of the systemcall is restricted to privileged users (Unix
\verb|root|) while other users need to have the \verb|CAP_PERFMON| capability.
The systemcall also registers a LSM hook which can be used by any security
module to provide custom access policies. Additionally, the subsystem register
the \verb|perf_event_paranoid| parameter, accessible in \textit{
  /proc/sys/kernel / perf\_event\_paranoid}, which can be set to provide
stricter or looser access checks, as described in the systemcall manual page.

\section{Others}

There are some vendor and use-case specific systems.

TODO: expand this section

\begin{itemize}

  \item audit
  \item systemtap
  \item LTTng
  \item ktap
  \item dtrace4linux
  \item DTrace
  \item sysdig

\end{itemize}

TODO: give an overview of these systems
